\chapter{2017 年试题（当年科目代码 819）}

\begin{tishi}
\textbf{关于科目代码：}原卷封面所印为「考试科目代码及名称：\textbf{819} 工程流体力学」，
与 2015 年相同，\underline{不是}现行 818 代码；本册统一按 818 编排，\textbf{请以原卷为准}。
不过本套试题的\textbf{内容全部落在现行 818 大纲范围内}——选择题、填空题、判断题、计算题中
均无明渠流、气体动力学、流体机械等旧大纲独有内容，复习时按现行大纲正常对待即可。

\textbf{试卷说明（原卷原文）：}试题附在答卷内交回，答案必须写在答题纸上，写在试题纸上无效。

\textbf{结构：}选择题 30 分（10 小题，每小题 3 分）、填空题 30 分（10 空，每空 3 分）、
判断题 30 分（10 小题，每小题 3 分）、计算题 60 分（6 小题，10、8、10、12、10、10 分）。满分 150 分。

\textbf{关于图示：}原卷图一律不作想象重绘，正文中标注「原卷图 $N$」的文字图示说明，
只是依据原卷写成的\emph{解题用文字描述}，\textbf{不等于原卷图形}；
凡描述不到之处均注明"原卷图略，请对照原卷"，影印件见本册附录。

\textbf{OCR 校订说明：}沿程阻力系数原卷 OCR 作「入」，按流体力学习惯改正为 $\lambda$；
泵效率 $\eta_1$、电动机效率 $\eta_2$ 原卷 OCR 作「ni」「n2」；水头损失单位 mH$_2$O 原卷 OCR 作「mH20」；
空气密度 $\rho=1.25\,\mathrm{kg/m^3}$ 原卷 OCR 作「p=1.25Kg/m3」。
以上均属字形错读的改正，\textbf{各题题设数值一律未作改动}。
\end{tishi}

\section{一、选择题（共 30 分，每小题 3 分）}

\begin{ti}
  \item 某点的绝对压强为 $65000$ Pa，当地大气压为 $0.1$ MPa，则该点的真空压强是（\quad）。\hfill\xstarfull{4}\yiju{E4 2018 选择 3（U 形水银测压计）；E4 2015 填空（U 形压差计）}\nandu{易}
  \fourchoices{$65000$ Pa}{$35000$ Pa}{$55000$ Pa}{$165000$ Pa}

  \item 并联管道 1、2，两管的直径相同，沿程阻力系数 $\lambda$ 相同，长度 $L_2=3L_1$，
  则通过流量 $Q_1$ 和 $Q_2$ 之间的关系式为（\quad）。\hfill\xstarfull{4}\yiju{E1 slide33；E6 教材 5.7}\nandu{中}
  \fourchoices{$Q_1=Q_2$}{$Q_1=1.5Q_2$}{$Q_1=1.73Q_2$}{不定}

  \item 圆管内的流动状态为层流时，沿程阻力与平均速度的（\quad）次方成正比。\hfill\xstarfull{5}\yiju{E4 2022 单选 10（管径减半压强损失 16 倍）；E4 2015 填空（层流平均速度＝0.5 倍最大速度）}\nandu{易}
  \fourchoices{0}{2}{3}{1}

  \item 空气的动力粘度随温度的升高而（\quad）。\hfill\xstarfull{5}\yiju{E4 2026 计算 1（牛顿流体斜板静压强）；E5 题库选择 2（汽油＝牛顿流体）}\nandu{易}
  \fourchoices{增大}{减小}{不变}{不定}

  \item 水平放置的渐扩管，如忽略水头损失，则上游断面形心压强 $p_1$ 和下游断面形心压强 $p_2$
  的大小关系为（\quad）。\hfill\xstarfull{5}\yiju{E1 slide16"计算题：伯努利方程"；E1 slide33 与 E2 slide4"第四章＝考试计算题重点"}\nandu{中}
  \fourchoices{$p_1>p_2$}{$p_1=p_2$}{$p_1<p_2$}{不定}

  \item 如图 1 所示，粘性流体在等直管道中流动，在过流断面 1、2 上有 $A$、$B$、$C$ 三点，
  则下面关系式成立的是（\quad）。\hfill\xstarfull{5}\yiju{E1 slide16"计算题：伯努利方程"；E1 slide33 与 E2 slide4"第四章＝考试计算题重点"}\nandu{难}
  \begin{choices}
    \item $z_A+\dfrac{p_A}{\gamma}=z_B+\dfrac{p_B}{\gamma}$
    \item $z_B+\dfrac{p_B}{\gamma}=z_C+\dfrac{p_C}{\gamma}$
    \item $z_A+\dfrac{p_A}{\gamma}=z_C+\dfrac{p_C}{\gamma}$
    \item $z_A+\dfrac{p_A}{\gamma}>z_C+\dfrac{p_C}{\gamma}$
  \end{choices}
  \tuyi{原卷图 1：一等直管道内画有数条水平流线（水流自左向右）；两条铅垂线分别为过流断面 1（上游）
  与 2（下游）；$A$、$C$ 两点均在上游断面 1 上，其中 $A$ 在上、$C$ 在下；$B$ 点在下游断面 2 上，
  并与 $A$ 位于同一条流线上。（原卷影印见本册附录）}

  \item 压力输水管同种流体的模型实验，已知长度比为 4，则两者的流量比为（\quad）。\hfill\xstarfull{3}\yiju{E4 2014 计算 8（溢流坝模型相似）；E5 题库选择 15}\nandu{中}
  \fourchoices{4}{8}{16}{32}

  \item 流量为 $Q$、速度为 $v$ 的射流冲击一块与流向垂直的平板，则平板受到的冲击力为（\quad）。\hfill\xstarfull{5}\yiju{E4 2015 计算（水平弯管求 $F_x,F_y$）；E4 2026 计算 3（教材习题 4.12）；E1 slide16}\nandu{易}
  \fourchoices{$Qv$}{$Qv^{2}$}{$\rho Qv$}{$\rho Qv^{2}$}

  \item 如图 2 所示的封闭容器上装有 U 型水银测压计，其中 1、2、3 点位于同一水平面上，
  其压强关系为（\quad）。\hfill\xstarfull{4}\yiju{E4 2018 选择 3（U 形水银测压计）；E4 2015 填空（U 形压差计）}\nandu{中}
  \fourchoices{$p_1>p_2>p_3$}{$p_1=p_2=p_3$}{$p_1<p_2<p_3$}{$p_2<p_1<p_3$}
  \tuyi{原卷图 2：一封闭容器（内盛水），液面上的压强为 $p_0$；容器侧壁接一根 U 型水银测压计，
  U 形管底部盛水银；1、2、3 三点位于同一水平面上——3 点在容器内水中，2 点在容器通向 U 形管的液体中，
  1 点在水银面上方的右侧管内（水柱中）。（原卷影印见本册附录）}

  \item 伯努利方程中 $z+\dfrac{p}{\rho g}+\dfrac{\alpha v^{2}}{2g}$ 表示（\quad）流体具有的机械能。\hfill\xstarfull{5}\yiju{E1 slide16"计算题：伯努利方程"；E1 slide33 与 E2 slide4"第四章＝考试计算题重点"}\nandu{易}
  \fourchoices{单位质量}{单位体积}{单位重量}{单位时间}
\end{ti}

\section{二、填空题（共 30 分，每空 3 分）}

\noindent 说明：物理量应注明单位。

\begin{ti}
  \item 流体微团的运动形式有 \underline{\hspace{2.5cm}} 、旋转和 \underline{\hspace{2.5cm}} 。\hfill\xstarfull{4}\yiju{E5 题库选择 13（涡量为零）；E6 教材 3.4}\nandu{易}

  \item 如图 3 所示的密封容器，当已知测压管高出液面 $h=1.5$ m 时，液面相对压强
  $p_0=$ \underline{\hspace{1.8cm}} Pa。容器盛的液体是汽油，其重度 $\gamma=7.35\,\mathrm{kN/m^3}$。\hfill\xstarfull{5}\yiju{E4 2022 单选 3（测压管水头线坡度）、单选 7（虹吸管）；E1 slide33"计算题基础入门"}\nandu{中}
  \tuyi{原卷图 3：一密封容器（内盛汽油），液面上方作用的压强为 $p_0$；容器侧壁接一根测压管，
  管内液面高出容器内液面 $h=1.5$ m。（原卷影印见本册附录）}

  \item 单位质量力的国际单位是 \underline{\hspace{3cm}} 。\hfill\xstarfull{3}\yiju{E4 2018 选择 1（作用于流体的质量力）；E5 题库选择 3（体积力有势）}\nandu{易}

  \item 圆管层流，断面平均流速为 $0.2$ m/s，则管轴处的切应力为 \underline{\hspace{2.5cm}} ，
  管轴上的流速为 \underline{\hspace{2.5cm}} 。\hfill\xstarfull{5}\yiju{E4 2022 单选 10（管径减半压强损失 16 倍）；E4 2015 填空（层流平均速度＝0.5 倍最大速度）}\nandu{中}

  \item 如图 4 所示的 $A$ 点的压强水头为 \underline{\hspace{2cm}} ，$A$ 点的测压管水头为 \underline{\hspace{2cm}} ，
  $B$ 点的压强水头为 \underline{\hspace{2cm}} ，$C$ 点的压强水头为 \underline{\hspace{2cm}} 。
  （$D$ 点闸门关闭，以 $D$ 点所在的水平面为基准面）\hfill\xstarfull{5}\yiju{E4 2022 单选 3（测压管水头线坡度）、单选 7（虹吸管）；E1 slide33"计算题基础入门"}\nandu{中}
  \tuyi{原卷图 4：$A$ 点位于容器上部的水面处；$B$ 点在 $A$ 点以下 $2$ m 处；$C$ 点在 $B$ 点以下
  $1$ m 的弯管转折处；自 $C$ 点接一段倾斜管通到末端 $D$ 点，$C$ 点与 $D$ 点的铅垂高差为 $3$ m；
  $D$ 点闸门关闭，并以 $D$ 点所在的水平面为基准面。（原卷影印见本册附录）}
\end{ti}

\section{三、判断题（正确的打"√"，错误的打"×"，共 30 分，每小题 3 分）}

\begin{ti}
  \item 在恒定流动中，流线与迹线重合。\hfill\xstarfull{5}\yiju{E4 2015 计算 1（流线方程、流动方向判断）}\nandu{易}

  \item 在均匀流动中，时变加速度为零。\hfill\xstarfull{5}\yiju{E4 各年选择/填空（质点导数与加速度）；E1 slide33"公式较多，重点记忆"}\nandu{中}

  \item 理想流体的粘度等于零。\hfill\xstarfull{3}\yiju{E4 2026 简答 2（理想流体）；E6 教材 1.3}\nandu{易}

  \item 雷诺数的物理意义为惯性力和粘滞力之比。\hfill\xstarfull{5}\yiju{E4 2022 单选 9（两圆管流态判别）；E4 2026 简答 3（层流与雷诺数关系）}\nandu{易}

  \item 无旋流动是指流线为直线的流动。\hfill\xstarfull{4}\yiju{E5 题库选择 13（涡量为零）；E6 教材 3.4}\nandu{中}

  \item 在均质连通的同种静止液体中，任一水平面上各点压强相同。\hfill\xstarfull{5}\yiju{E4 2022 单选 3（测压管水头线坡度）、单选 7（虹吸管）；E1 slide33"计算题基础入门"}\nandu{易}

  \item 圆管湍流粗糙区的沿程阻力系数 $\lambda$ 与管壁相对粗糙度 $\Delta/d$ 有关。\hfill\xstarfull{3}\yiju{E6 教材 5.4；E2 slide5}\nandu{中}

  \item 动力粘度 $\mu$ 的量纲为 $\mathrm{MLT^{-1}}$。\hfill\xstarfull{5}\yiju{E4 2026 计算 1（牛顿流体斜板静压强）；E5 题库选择 2（汽油＝牛顿流体）}\nandu{中}

  \item 粘性流体的测压管水头线沿程一定是下降的。\hfill\xstarfull{5}\yiju{E1 slide16"计算题：伯努利方程"；E1 slide33 与 E2 slide4"第四章＝考试计算题重点"}\nandu{中}

  \item 当流体处于静止状态时，流体的切应力不会产生。\hfill\xstarfull{4}\yiju{E2 slide4；E5 概念题高频}\nandu{易}
\end{ti}

\section{四、计算题（共 60 分）}

\begin{ti}
  \item 如图 5 所示，已知离心泵的提水高度 $z=20$ m，抽水流量 $Q=35$ L/s，泵的效率 $\eta_1=0.82$。
  若吸水管路和压水管路总水头损失 $h_w=1.5$ mH$_2$O，电动机的效率 $\eta_2=0.95$，
  试求：电动机的功率 $P$。（10 分）\hfill\xstarfull{4}\yiju{E5 题库计算 4（水泵工况点与节流调节）；E6 教材 4.2}\nandu{中}
  \tuyi{原卷图 5：左下方为吸水池，池中自由液面为断面 1--1；离心泵经吸水管自池中吸水，
  再经压水管将水送至右上方出水池，池中自由液面为断面 2--2；图上标出两自由液面的高差 $z$。
  （原卷影印见本册附录）}

  \item 汽车以 $80$ km/h 的时速行驶，其迎风面积为 $A=2\,\mathrm{m^2}$，阻力系数为 $C_D=0.4$，
  空气的密度为 $\rho=1.25\,\mathrm{kg/m^3}$，试求汽车克服空气阻力所消耗的功率。（8 分）\hfill\xstarfull{4}\yiju{E4 2014 选择 10；E4 2016 填空 4；E4 2019 判断 2；E4 2022 单选 9；E5 题库选择 12}\nandu{中}

  \item 如图 6 所示，水从直径 $d$、长 $L$ 的铅垂管路流入大气中，水箱中液面高度为 $h$，
  管路局部阻力可忽略，沿程阻力系数为 $\lambda$。（10 分）\hfill\xstarfull{5}\yiju{E1 slide16"计算题：伯努利方程"；E1 slide33 与 E2 slide4"第四章＝考试计算题重点"}\nandu{中}
  \begin{xiaowen}
    \item 求管路起始断面 $A$ 处压强。
    \item $h$ 等于多少时，可使 $A$ 点的压强等于大气压。
  \end{xiaowen}
  \tuyi{原卷图 6：水箱中的水经一根铅垂管路流入大气，水箱中液面高度为 $h$；管路直径 $d$、长 $L$，
  $A$ 为管路起始断面（水箱底部出口处），液面到 $A$ 断面的高度即 $h$。（原卷影印见本册附录）}

  \item 如图 7 所示，水流经水平弯管流入大气，已知 $d_1=100$ mm，$d_2=75$ mm，
  $v_1=1.5$ m/s，$\theta=30^\circ$。若不计水头损失，试求水流对弯管的作用力 $F_x$、$F_y$。（12 分）\hfill\xstarfull{5}\yiju{E4 2015 计算（水平弯管求 $F_x,F_y$）；E4 2026 计算 3（教材习题 4.12）；E1 slide16}\nandu{难}
  \tuyi{原卷图 7：水平放置的渐变弯管，进口断面 1--1 直径 $d_1=100$ mm、流速 $v_1=1.5$ m/s 沿水平方向；
  出口断面 2--2 直径 $d_2=75$ mm，出口轴线与进口轴线成 $\theta=30^\circ$ 夹角，出口水流自由流入大气。
  （原卷影印见本册附录）}

  \item 长管输送液体只计沿程损失。如图 8 所示，当 $H$、$L$ 一定、沿程损失为 $H/3$ 时
  管路输送的功率为最大；已知 $H=127.4$ m，$L=500$ m，管路末端可用水头 $h=2H/3$，
  管路末端可用功率为 $1000$ kW，$\lambda=0.024$，求管路的输送流量与管路直径。（10 分）\hfill\xstarfull{5}\yiju{E5 题库计算 2（水箱垂直管道出流 $Q$ 与管长 $l$）；E1 slide33}\nandu{难}
  \tuyi{原卷图 8：左侧为上游水池及其自由液面，液面至管路末端基准面的总水头为 $H$；
  管路水平、长为 $L$；图中自左向右下降的直线为管路的总水头线，其下降量即沿程水头损失 $h_f$；
  管路末端剩余的水头为 $h$（$H=h_f+h$），末端所画圆圈示意利用 $h$ 做功的装置；
  图中未标注管径。（原卷影印见本册附录）}

  \item 如图 9 所示，一水坝受水面为一抛物线，顶点在 $O$ 点，水深 $h=50$ m 处抛物线到
  抛物线对称轴的距离为 $12.5$ m，求水作用于单位宽度坝体的水平分力和垂直分力。（10 分）\hfill\xstarfull{5}\yiju{E5 题库计算 1（弧形闸门 $F_x,F_z$ 与力矩）；E1 slide33}\nandu{难}
  \tuyi{原卷图 9：坝体迎水面为一条抛物线，抛物线顶点为 $O$ 点，位于水深 $h=50$ m 处；
  过 $O$ 点的铅垂点划线为抛物线的对称轴；水面线与抛物线的交点（坝顶）$A$ 在对称轴右侧，
  其到对称轴的水平距离为 $12.5$ m；水位于抛物线左侧。（原卷影印见本册附录）}
\end{ti}
